\clearpage
\chapter*{Errata da exposição n\textsuperscript{o} 221}
\addcontentsline{toc}{chapter}{Errata da exposição n\textsuperscript{o} 221}
\setcounter{footnote}{3}
\src{302}

\noindent\textbf{6. TDTE IV: Os esquemas de Hilbert.}

\begin{center}
[Seminário Bourbaki, t. 13, 1960/61, n\textsuperscript{o} 221, 28 p.]
\end{center}

\noindent\textbf{Página 6, teorema 2.1, 3\textsuperscript{a} linha.} --- Em vez de
«é necessário», leia-se «é necessário e suficiente».

\noindent\textbf{Página 15, linha 13.} --- Em vez de «$P\neq Q$», leia-se
«$P<Q$».

\noindent\textbf{Linha 15.} --- Em vez de «Se, ao contrário, $P=Q$», leia-se
«No caso contrário, teremos $P(n)\geq Q(n)$ para $n$ grande».

\noindent\textbf{Linha 18, fórmula (*).} --- Em vez de «$<$» e «$>$»,
leia-se «$\leq$» e «$\geq$».

\noindent\textbf{5.ª linha a partir de baixo.} --- Em vez de «$r$», leia-se
«$Z$».

\noindent\textbf{Página 18, observação 3.9.} --- Assinalemos que o estudo das
componentes conexas dos esquemas de Hilbert sobre um corpo algebricamente
fechado foi realizado por HARTSHORNE. Esse autor demonstra que os
$\operatorname{Hilb}^{P}_{r}$ são conexos e determina os pares $(r,P)$ para
os quais $\operatorname{Hilb}^{P}_{r}\neq\varnothing$\footnote{HARTSHORNE
(R.). --- Connectedness of the Hilbert scheme (Tese, Harvard, 1962).}.

\noindent\textbf{Página 23, 5.ª linha a partir de baixo.} --- Em vez de «uma
componente», leia-se «de uma componente».

\noindent\textbf{Páginas 23 e 24, observação 5.5.} --- A propósito do exemplo
de ZAPPA, assinalemos que MUMFORD acaba até mesmo de construir uma componente
irredutível do esquema de Hilbert para $\mathbf P^3_{\mathbf C}$ (cujos pontos
gerais representam curvas não singulares de grau 14 e gênero 24), que não é
reduzida em seu ponto genérico. Realizando o blow-up de uma das curvas obtidas,
ele obtém também um esquema projetivo regular de dimensão 3 sobre $\mathbf C$,
cujo esquema formal de moduli não é reduzido em seu ponto genérico ou, o que
equivale ao mesmo, cuja variedade local de moduli, no sentido da geometria
analítica, sobre $\mathbf C$ (cf. Seminário Cartan, t. 13, 1960/61), não é
reduzida em nenhum de seus pontos.
